\begin{proposition}[Integral homology and monodromy of the smooth fibre]\hypertarget{tag-02ZS}{}\label{push:src:A.4.2}\label{push:A42}
\StatementRecord{push:src:A.4.2}{proved}{02ZS}{\SeriesRef{push:src:A.2.2}, \SeriesRef{push:src:A.4.1}}

The untwisted smooth fibre has torsion-free integral homology
with Betti numbers
\[
 (b_0,b_1,b_2,b_3,b_4,b_5,b_6)=(1,1,8,14,8,1,1).
\]
The positive-$u$ monodromy is the identity outside degree three.
In degree three there is an integral exact sequence
\[
 0\longrightarrow \mathcal L/\langle F,v,m\rangle
 \longrightarrow H_3(X(u);\mathbb Z)
 \longrightarrow K\longrightarrow0,\qquad
 K=\{a\in\mathcal L:aF=av=am=0\},
\]
in which $T_u-I$ is zero on the left term and, on the quotient,
is the map
\[
 K\longrightarrow \mathcal L/\langle F,v,m\rangle,\qquad
 a\longmapsto a\bmod\langle F,v,m\rangle.
\]
The sign convention orients the neck circle by increasing
$\arg x$ in $xy=u$; the boundary convention for a transported
chain is chosen so that a positive $u$ loop adds this circle
times its boundary cycle. With these conventions
$(T_u-I)^2=0$, its rank is six, and its nonzero Smith factors
are $1,1,1,1,1,11$.
\end{proposition}


\begin{proof}
\begin{enumerate}[label=\arabic*.,leftmargin=*]
\item We use the semistable model \ref{push:src:A.2.2}. Remove disjoint
open tubular neighbourhoods of the two conductor sections
from $R^+$, and denote the remaining compact manifold by
$Z$. Its two boundary circle bundles are identified by the
plumbing; let $E$ denote the first. The Euler class of
$E\to S$ is $n=m-F$. It is primitive. The Thom sequence
for $(R^+,Z)$ gives
\[
 H_5(Z)=\mathbb Z,\quad H_4(Z)=0,\quad
 H_3(Z)=\mathcal L/\langle F,m\rangle\simeq\mathbb Z^8,
 \quad H_2(Z)\simeq\mathbb Z^{10},\quad H_1(Z)=0.
\]
We spell out the rank and saturation assertions. If $H$ is
the projective-bundle class restricting to zero on $s_0$
and to $-m$ on $s_\infty$, and $A$ is the exceptional
divisor class, the two restriction maps on $H^2(R^+)$ are
\[
 (\alpha,h,a)\longmapsto\alpha+aF,\qquad
 (\alpha,h,a)\longmapsto\alpha-hm.
\]
Their combined image is primitive of rank twelve, and its
cokernel is $\mathcal L/\langle F,m\rangle$. In degree two,
intersection with the two sections is onto $\mathbb Z^2$:
a generic projective fibre meets both once, and an exceptional
projective fibre meets just the section containing the centre
once. In degree six the map is the diagonal
$\mathbb Z\to\mathbb Z^2$. These are precisely the Thom
maps giving the asserted groups.

\item The Gysin sequence of the primitive circle bundle gives
\[
 H_2(E)=n^\perp\subset\mathcal L,\qquad
 H_3(E)=\mathcal L/\mathbb Z n,\qquad
 H_1(E)=H_4(E)=0,\quad H_5(E)=\mathbb Z.
\]
The smooth fibre is obtained from $Z$ by the boundary
mapping cylinder. In its Mayer--Vietoris sequence the
degree-two difference map is the normalization difference
$d_2^+$ restricted to $n^\perp$. In the coefficient
coordinates $(s,t,h)=(aF,av,am)$ the constraint is $h=s$,
so this map has primitive rank two and kernel $K$.
The degree-three difference map is, modulo $F,m$,
$(T_*-I)a=(aF)v$, hence has primitive image $\mathbb Z v$.
These statements use the surjectivity of the three
intersection coordinates proved in \ref{push:src:A.4.1}, so no integral
index is omitted. The degree-five difference is zero:
the plumbing reverses the normal circle orientation, and
the two boundary fundamental classes sum to zero in $Z$.
The degree-zero difference is zero as well. Exactness now
gives the claimed groups and the displayed degree-three
extension. Poincar\'e duality, or the same integral exact
sequence, confirms the absence of torsion in complementary
degrees.

\item We also compute monodromy directly from this decomposition.
In the neck $xy=u$, a positive circuit of $u$ rotates one
boundary normal coordinate once relative to the other.
A cycle in $Z$ is unchanged. A three-cycle whose cut-open
boundary represents $a\in K$ acquires the tube obtained
by sweeping $a$ once around that normal circle. The Gysin
description sends this tube to the class of $a$ in
$H_3(E)=\mathcal L/\mathbb Z n$, and then to its class in
$\mathcal L/\langle F,m,v\rangle$. This is exactly the
map in the statement. The same description in the other
degrees is zero, using the groups and difference maps just
computed. A tube has no cut-open boundary, which proves
$(T_u-I)^2=0$ without an appeal to a monodromy theorem.

\item Finally $K\cap\langle F,m,v\rangle=\mathbb ZF$, since
the Gram matrix of $m-F,v$ has determinant eleven.
Hence the rank is six. In the unimodular $E_8$ summand,
the primitive rank-two lattice spanned by $m-F,v$ has
discriminant eleven; its orthogonal complement maps into
the dual quotient with cokernel of order eleven. Thus the nonzero Smith factors are five ones and eleven.
\end{enumerate}
\end{proof}
